\section{Discusión de Resultados}\label{sec:discussion}

\subsection{Idoneidad del algoritmo frente a las propiedades de los datos}
La elección de un ensamble Random Forest no fue arbitraria, sino una consecuencia
directa de las propiedades intrínsecas del conjunto de datos reveladas durante el
análisis exploratorio. Cuatro características del dominio favorecieron
específicamente a este algoritmo frente a alternativas.

Primero, el EDA evidenció \emph{alta multicolinealidad} entre grupos de sensores
relacionados (variantes de viento y de temperatura). Los árboles de decisión que
componen el bosque realizan particiones ortogonales del espacio de
características, lo que los hace intrínsecamente robustos a la redundancia lineal
entre predictores; combinados con la reducción de dimensionalidad vía PCA, que
descorrelaciona explícitamente el espacio, se obtuvo una sinergia geométrica
favorable a las particiones axiales del ensamble. Un modelo lineal como la
regresión logística habría sufrido de inestabilidad en sus coeficientes ante esta
multicolinealidad.

Segundo, el conjunto presentaba \emph{ruido de instrumentación} en forma del
código \texttt{-999} y \emph{outliers univariados legítimos} (precipitación de
hasta $470$\,mm). Random Forest es notablemente tolerante a valores atípicos, dado
que sus divisiones se basan en umbrales ordinales y no en distancias métricas
absolutas; esto contrasta con algoritmos sensibles a la escala como SVM con
kernel RBF o $k$-NN, cuyo desempeño se habría degradado sin un tratamiento de
outliers más agresivo. La combinación con \texttt{RobustScaler} reforzó esta
tolerancia al centrar y escalar mediante estadísticos resistentes (mediana e
IQR).

Tercero, el problema exhibió un \emph{desbalance de clases severo} ($16.24\%$
frente a $83.76\%$). La capacidad de Random Forest de incorporar ponderación de
clases (\texttt{class\_weight='balanced'}) permitió reajustar la función de
decisión hacia la clase minoritaria sin recurrir a generación sintética de
muestras, cuyo uso ---según se discutió en el estado del arte--- puede magnificar
la complejidad de datos ya ruidosos. La optimización sobre F1-Score Macro alineó
directamente la función objetivo con la métrica operativamente relevante.

Cuarto, la señal predictiva resultó \emph{no trivial en su ubicación espectral}:
el análisis de impureza de Gini reveló que la capacidad discriminante reside en
componentes de varianza intermedia (PC3, PC2) y no en la dirección dominante
(PC1). La naturaleza no paramétrica de Random Forest, capaz de capturar
interacciones no lineales entre componentes sin asumir una forma funcional
predeterminada, resultó adecuada para explotar una señal distribuida de esta
manera, algo que un clasificador lineal sobre PC1 no habría capturado.

\subsection{Interpretación del compromiso Precision--Recall}
El resultado central ---un \emph{Recall} de $92.34\%$ frente a una \emph{Precision}
de $61.09\%$ en la clase de helada--- debe interpretarse a la luz del contexto de
aplicación y no como una deficiencia. En un sistema de alerta temprana
agrometeorológica, la matriz de costos es marcadamente asimétrica: el costo de un
falso negativo (una helada no anticipada que destruye un cultivo) supera
ampliamente el de un falso positivo (una alerta preventiva que no se materializa).
La ponderación de clases y la optimización sobre F1-Macro desplazaron
deliberadamente el modelo hacia la sensibilidad, priorizando la detección
exhaustiva de eventos críticos. El elevado ROC-AUC ($0.9690$) confirma que esta
sensibilidad no se obtuvo a costa de una capacidad discriminante global pobre.

\subsection{Amenazas a la validez interna}
La validez interna se refiere a si las conclusiones del experimento son correctas
dentro del entorno controlado del estudio. Se identifican las siguientes
amenazas.

\emph{Prevención de fuga de información.} El riesgo más severo ---la fuga por
predicción tautológica--- fue mitigado explícitamente al eliminar
\texttt{temp\_min\_2m} del espacio de predictores tras derivar de ella la etiqueta.
Asimismo, todos los parámetros de imputación y escalamiento se ajustaron
exclusivamente sobre el conjunto de entrenamiento, transformando el de prueba de
forma ciega. No obstante, el ajuste de PCA también se realizó sobre el conjunto
de entrenamiento, lo cual es metodológicamente correcto, pero su interacción con
la ponderación de clases no fue objeto de una validación anidada, lo que
constituye una amenaza residual menor.

\emph{Sesgo de selección en la matriz de correlación.} El análisis de correlación
de Pearson se computó sobre una submuestra de $20\,000$ registros completos, no
sobre la totalidad del conjunto. Aunque la submuestra fue aleatoria y con semilla
fija, la exclusión de registros con nulos (\texttt{dropna}) podría introducir un
sesgo si la ausencia de datos no fuese completamente aleatoria (MCAR). Dado que
los faltantes provienen de fallos de sensor uniformemente distribuidos
($0.03\%$), este sesgo se considera despreciable, pero no puede descartarse
formalmente.

\emph{Ausencia de validación estadística.} Como se detalló en la
Sección~\ref{sec:results}, el experimento no incluye múltiples ejecuciones
independientes ni modelos de línea base, por lo que no es posible afirmar con
respaldo estadístico que el desempeño observado sea significativamente superior al
de alternativas. Las conclusiones se sustentan en un único punto experimental, lo
que constituye la principal amenaza a la validez interna del estudio.

\subsection{Amenazas a la validez externa}
La validez externa concierne a la capacidad de generalización del modelo a
entornos productivos fuera del conjunto de prueba.

\emph{Representatividad geográfica y temporal.} El conjunto abarca estaciones
peruanas en un rango de latitud $[-18.01, -8.19]$ y altitud
$[1079, 4691]$\,m.s.n.m., a lo largo del periodo 2000--2025. La generalización a
microclimas no representados en esta cobertura, o a regiones de menor altitud, no
está garantizada. La helada es un fenómeno intrínsecamente local, y un modelo
entrenado sobre estas estaciones podría degradar su desempeño en ubicaciones con
regímenes térmicos distintos.

\emph{Deriva de distribución (\emph{dataset shift}).} En un despliegue productivo,
los patrones climáticos pueden desviarse de la distribución de entrenamiento por
efecto de la variabilidad climática de largo plazo. Un modelo estático como el
presentado requeriría reentrenamiento periódico para mantener su validez ante
posibles cambios en las distribuciones marginales de las variables.

\emph{Dependencia de la calidad de instrumentación.} El pipeline asume que los
fallos de sensor se codifican mediante el valor \texttt{-999}. En un entorno
productivo con sensores heterogéneos que empleen otros códigos de error o que
produzcan lecturas corruptas de forma no señalizada, el mecanismo de corrección
podría no capturar todas las anomalías, comprometiendo la robustez del sistema.

\emph{Ausencia de estructura temporal explícita.} El modelo trata cada
observación de forma independiente, sin incorporar descriptores espaciotemporales
(estacionalidad, tendencia). En producción, la predicción anticipada de heladas
se beneficiaría de una formulación que capture la dinámica temporal, lo que
sugiere que el desempeño en un escenario de pronóstico real (a $t+k$ horas) podría
diferir del obtenido en la clasificación sobre observaciones concurrentes.